
\documentclass{article}
\usepackage{colortbl}
\usepackage{makecell}
\usepackage{multirow}
\usepackage{supertabular}

\begin{document}

\newcounter{utterance}

\centering \large Interaction Transcript for game `clean\_up', experiment `0\_easy\_7obj\_cs', episode 1 with gemma{-}4{-}26b{-}a4b{-}it.
\vspace{24pt}

{ \footnotesize  \setcounter{utterance}{1}
\setlength{\tabcolsep}{0pt}
\begin{supertabular}{c@{$\;$}|p{.15\linewidth}@{}p{.15\linewidth}p{.15\linewidth}p{.15\linewidth}p{.15\linewidth}p{.15\linewidth}}
    \# & $\;$A & \multicolumn{4}{c}{Game Master} & $\;\:$B\\
    \hline

    \theutterance \stepcounter{utterance}  
    & & \multicolumn{4}{p{0.6\linewidth}}{
        \cellcolor[rgb]{0.9,0.9,0.9}{
            \makecell[{{p{\linewidth}}}]{
                \texttt{\tiny{[P1$\langle$GM]}}
                \texttt{Jsem váš herní mistr a hrajete společnou hru s následující mřížkou jako herním plánem:} \\
\\ 
\texttt{```} \\
\texttt{ 1234567} \\
\texttt{╔═══════╗} \\
\texttt{║◌◌◌◌◌T◌║ 1} \\
\texttt{║P◌◌◌◌◌◌║ 2} \\
\texttt{║◌◌AD◌┌─╢ 3} \\
\texttt{║◌◌◌┌─┤S║ 4} \\
\texttt{║O◌◌│◌├─╢ 5} \\
\texttt{╟───┼─┘◌║ 6} \\
\texttt{║◌M◌│◌◌◌║ 7} \\
\texttt{╚═══╧═══╝} \\
\\ 
\texttt{```} \\
\\ 
\texttt{* Horní okraj zobrazuje souřadnice x rostoucí doprava a pravý okraj souřadnice y rostoucí dolů.} \\
\texttt{* Na vaší mřížce jsou náhodně umístěny následující objekty: 'P', 'O', 'T', 'S', 'D', 'A', 'M'.} \\
\\ 
\\ 
\texttt{Druhý hráč vidí variantu herní desky, kde jsou objekty umístěny na různých náhodných místech. Vy nevidíte desku druhého hráče a on nevidí tu vaši.} \\
\\ 
\\ 
\texttt{**Cíl:**} \\
\\ 
\texttt{Oba hráči musí přesunout objekty na svém pozadí tak, aby stejné objekty skončily na stejných souřadnicích. Musíte komunikovat s druhým hráčem, abyste se dohodli na společné cílové konfiguraci.} \\
\\ 
\\ 
\texttt{**Pravidla:**} \\
\\ 
\texttt{* V každém tahu můžete poslat přesně jeden z následujících dvou příkazů:} \\
\texttt{ 1. `ŘEKNI: <ZPRÁVA>`: odeslání zprávy (vše až do dalšího zalomení řádku) druhému hráči. Přepošlu ji vašemu partnerovi.} \\
\texttt{2. `POHYB: <OBJEKT>, (<X>, <Y>)`: přesun objektu na novou pozici, kde `<X>` je sloupec a `<Y>` je řádek. Budu vás informovat, zda byl váš přesun úspěšný, nebo ne.} \\
\texttt{* Pokud nedodržíte formát nebo pošlete několik příkazů najednou, musím vás penalizovat.} \\
\texttt{* Pokud oba hráči nasbírají více než 16 penalizací, oba prohrajete hru.} \\
\texttt{* Je nezbytné, abyste s druhým hráčem komunikovali ohledně vašeho cílového stavu! *Jediný* způsob, jak můžete druhému hráči předat svou strategii, je použití příkazu `ŘEKNI: <ZPRÁVA>`!} \\
\\ 
\\ 
\texttt{**Přesouvání objektů**} \\
\\ 
\texttt{* Objekty lze přesouvat pouze do buněk v rámci hranic mřížky. Cílová buňka musí být prázdná, tj. musí obsahovat pouze symbol '◌'.} \\
\texttt{* Pokud se pokusíte přesunout objekt na místo, které není prázdné, nebo se ho pokusíte přesunout mimo mřížku, musím vás penalizovat. Dostanete další pokus.} \\
\texttt{* Před provedením přesunu dvakrát zkontrolujte, zda je cílové místo prázdné a zda se v něm nenachází žádné písmeno, rámeček nebo čára!} \\
\\ 
\\ 
\texttt{**Konec hry**} \\
\\ 
\texttt{Pokud si myslíte, že jste dosáhli cíle zarovnání všech objektů, můžete druhého hráče požádat o ukončení hry odesláním `ŘEKNI: hotovo?`. Pokud vás druhý hráč požádá o ukončení hry a vy odpovíte `ŘEKNI: hotovo!`, hra skončí.} \\
\\ 
\texttt{Oba hráči vyhrají, pokud hra skončí během 28 kol, přičemž jedno kolo je definováno jako odeslání platného příkazu oběma hráči.} \\
\\ 
\texttt{**Bodování:**} \\
\\ 
\texttt{Čím blíže jsou stejné objekty na obou herních deskách, tím více bodů získáte. Penalizace vám body snižují. Podaří se vám překonat rekord?} \\
\\ 
\\ 
\texttt{Pošlete prosím zprávu druhému hráči a začněte hru!} \\
            }
        }
    }
    & \\ \\

    \theutterance \stepcounter{utterance}  
    & \multicolumn{4}{p{0.6\linewidth}}{
        \cellcolor[rgb]{0.9,0.9,0.9}{
            \makecell[{{p{\linewidth}}}]{
                \texttt{\tiny{[P1$\rangle$GM]}}
                \texttt{ŘEKNI: Ahoj! Musíme najít společnou konfiguraci. Já vidím objekty na těchto pozicích: P(1,2), O(1,5), T(6,1), S(7,4), A(4,3), D(3,3), M(2,7). Navrhuji, abychom se pokusili posunout všechny objekty do řady 1 a 2, aby to bylo přehledné. Co ty na to? Jaké máš ty?} \\
            }
        }
    }
    & & \\ \\

    \theutterance \stepcounter{utterance}  
    & & & \multicolumn{2}{p{0.3\linewidth}}{
        \cellcolor[rgb]{0.9,0.9,0.9}{
            \makecell[{{p{\linewidth}}}]{
                \texttt{\tiny{[GM$|$GM]}}
                \texttt{ŘEKNI: Ahoj! Musíme najít společnou konfiguraci. Já vidím objekty na těchto pozicích: P(1,2), O(1,5), T(6,1), S(7,4), A(4,3), D(3,3), M(2,7). Navrhuji, abychom se pokusili posunout všechny objekty do řady 1 a 2, aby to bylo přehledné. Co ty na to? Jaké máš ty?} \\
            }
        }
    }
    & & \\ \\

    \theutterance \stepcounter{utterance}  
    & & \multicolumn{4}{p{0.6\linewidth}}{
        \cellcolor[rgb]{0.9,0.9,0.9}{
            \makecell[{{p{\linewidth}}}]{
                \texttt{\tiny{[GM$\rangle$P2]}}
                \texttt{Jsem váš herní mistr a hrajete společnou hru s následující mřížkou jako herním plánem:} \\
\\ 
\texttt{```} \\
\texttt{ 1234567} \\
\texttt{╔═══════╗} \\
\texttt{║◌◌◌◌◌S◌║ 1} \\
\texttt{║◌◌◌◌O◌◌║ 2} \\
\texttt{║◌◌◌M◌┌─╢ 3} \\
\texttt{║◌◌A┌─┤◌║ 4} \\
\texttt{║◌◌◌│◌├─╢ 5} \\
\texttt{╟───┼─┘◌║ 6} \\
\texttt{║D◌T│P◌◌║ 7} \\
\texttt{╚═══╧═══╝} \\
\\ 
\texttt{```} \\
\\ 
\texttt{* Horní okraj zobrazuje souřadnice x rostoucí doprava a pravý okraj souřadnice y rostoucí dolů.} \\
\texttt{* Na vaší mřížce jsou náhodně umístěny následující objekty: 'P', 'O', 'T', 'S', 'D', 'A', 'M'.} \\
\\ 
\\ 
\texttt{Druhý hráč vidí variantu herní desky, kde jsou objekty umístěny na různých náhodných místech. Vy nevidíte desku druhého hráče a on nevidí tu vaši.} \\
\\ 
\\ 
\texttt{**Cíl:**} \\
\\ 
\texttt{Oba hráči musí přesunout objekty na svém pozadí tak, aby stejné objekty skončily na stejných souřadnicích. Musíte komunikovat s druhým hráčem, abyste se dohodli na společné cílové konfiguraci.} \\
\\ 
\\ 
\texttt{**Pravidla:**} \\
\\ 
\texttt{* V každém tahu můžete poslat přesně jeden z následujících dvou příkazů:} \\
\texttt{ 1. `ŘEKNI: <ZPRÁVA>`: odeslání zprávy (vše až do dalšího zalomení řádku) druhému hráči. Přepošlu ji vašemu partnerovi.} \\
\texttt{2. `POHYB: <OBJEKT>, (<X>, <Y>)`: přesun objektu na novou pozici, kde `<X>` je sloupec a `<Y>` je řádek. Budu vás informovat, zda byl váš přesun úspěšný, nebo ne.} \\
\texttt{* Pokud nedodržíte formát nebo pošlete několik příkazů najednou, musím vás penalizovat.} \\
\texttt{* Pokud oba hráči nasbírají více než 16 penalizací, oba prohrajete hru.} \\
\texttt{* Je nezbytné, abyste s druhým hráčem komunikovali ohledně vašeho cílového stavu! *Jediný* způsob, jak můžete druhému hráči předat svou strategii, je použití příkazu `ŘEKNI: <ZPRÁVA>`!} \\
\\ 
\\ 
\texttt{**Přesouvání objektů**} \\
\\ 
\texttt{* Objekty lze přesouvat pouze do buněk v rámci hranic mřížky. Cílová buňka musí být prázdná, tj. musí obsahovat pouze symbol '◌'.} \\
\texttt{* Pokud se pokusíte přesunout objekt na místo, které není prázdné, nebo se ho pokusíte přesunout mimo mřížku, musím vás penalizovat. Dostanete další pokus.} \\
\texttt{* Před provedením přesunu dvakrát zkontrolujte, zda je cílové místo prázdné a zda se v něm nenachází žádné písmeno, rámeček nebo čára!} \\
\\ 
\\ 
\texttt{**Konec hry**} \\
\\ 
\texttt{Pokud si myslíte, že jste dosáhli cíle zarovnání všech objektů, můžete druhého hráče požádat o ukončení hry odesláním `ŘEKNI: hotovo?`. Pokud vás druhý hráč požádá o ukončení hry a vy odpovíte `ŘEKNI: hotovo!`, hra skončí.} \\
\\ 
\texttt{Oba hráči vyhrají, pokud hra skončí během 28 kol, přičemž jedno kolo je definováno jako odeslání platného příkazu oběma hráči.} \\
\\ 
\texttt{**Bodování:**} \\
\\ 
\texttt{Čím blíže jsou stejné objekty na obou herních deskách, tím více bodů získáte. Penalizace vám body snižují. Podaří se vám překonat rekord?} \\
\\ 
\texttt{Druhý hráč zahájil hru odesláním této zprávy:} \\
\\ 
\texttt{"Ahoj! Musíme najít společnou konfiguraci. Já vidím objekty na těchto pozicích: P(1,2), O(1,5), T(6,1), S(7,4), A(4,3), D(3,3), M(2,7). Navrhuji, abychom se pokusili posunout všechny objekty do řady 1 a 2, aby to bylo přehledné. Co ty na to? Jaké máš ty?"} \\
\\ 
\texttt{Jaký je váš první příkaz?} \\
            }
        }
    }
    & \\ \\

    \theutterance \stepcounter{utterance}  
    & & & \multicolumn{4}{p{0.6\linewidth}}{
        \cellcolor[rgb]{0.9,0.9,0.9}{
            \makecell[{{p{\linewidth}}}]{
                \texttt{\tiny{[GM$\langle$P2]}}
                \texttt{ŘEKNI: Ahoj! Moje pozice jsou: S(6,1), O(5,2), M(4,3), A(3,4), D(1,7), T(3,7), P(5,7). Tvůj návrh na řady 1 a 2 je dobrý, ale tam je málo míst pro všech 7 objektů. Navrhuji raději vytvořit sloupec nebo dvě řady, například v horní části mřížky. Co kdybychom se pokusili o konfiguraci v prvních dvou řadách takto: S(1,1), O(2,1), T(3,1), P(4,1), A(5,1), M(6,1), D(7,1)? To by bylo v jedné řadě. Co si o tom myslíš?} \\
            }
        }
    }
     \\ \\

    \theutterance \stepcounter{utterance}  
    & & & \multicolumn{2}{p{0.3\linewidth}}{
        \cellcolor[rgb]{0.9,0.9,0.9}{
            \makecell[{{p{\linewidth}}}]{
                \texttt{\tiny{[GM$|$GM]}}
                \texttt{ŘEKNI: Ahoj! Moje pozice jsou: S(6,1), O(5,2), M(4,3), A(3,4), D(1,7), T(3,7), P(5,7). Tvůj návrh na řady 1 a 2 je dobrý, ale tam je málo míst pro všech 7 objektů. Navrhuji raději vytvořit sloupec nebo dvě řady, například v horní části mřížky. Co kdybychom se pokusili o konfiguraci v prvních dvou řadách takto: S(1,1), O(2,1), T(3,1), P(4,1), A(5,1), M(6,1), D(7,1)? To by bylo v jedné řadě. Co si o tom myslíš?} \\
            }
        }
    }
    & & \\ \\

    \theutterance \stepcounter{utterance}  
    & & \multicolumn{4}{p{0.6\linewidth}}{
        \cellcolor[rgb]{0.9,0.9,0.9}{
            \makecell[{{p{\linewidth}}}]{
                \texttt{\tiny{[P1$\langle$GM]}}
                \texttt{Vaše zpráva byla předána druhému hráči.} \\
\\ 
\\ 
\texttt{Právě hrajete kolo 1 z maximálně 28.} \\
\texttt{Dohromady jste nasbírali 0 z 16 pokut.} \\
\texttt{Druhý hráč poslal tuto zprávu:} \\
\\ 
\\ 
\texttt{"Ahoj! Moje pozice jsou: S(6,1), O(5,2), M(4,3), A(3,4), D(1,7), T(3,7), P(5,7). Tvůj návrh na řady 1 a 2 je dobrý, ale tam je málo míst pro všech 7 objektů. Navrhuji raději vytvořit sloupec nebo dvě řady, například v horní části mřížky. Co kdybychom se pokusili o konfiguraci v prvních dvou řadách takto: S(1,1), O(2,1), T(3,1), P(4,1), A(5,1), M(6,1), D(7,1)? To by bylo v jedné řadě. Co si o tom myslíš?"} \\
\\ 
\\ 
\texttt{Jaký je váš další příkaz?} \\
            }
        }
    }
    & \\ \\

    \theutterance \stepcounter{utterance}  
    & \multicolumn{4}{p{0.6\linewidth}}{
        \cellcolor[rgb]{0.9,0.9,0.9}{
            \makecell[{{p{\linewidth}}}]{
                \texttt{\tiny{[P1$\rangle$GM]}}
                \texttt{ŘEKNI: To je skvělý nápad! Jedna řada je velmi jednoduchá na koordinaci. Souhlasím s tvou konfigurací: S(1,1), O(2,1), T(3,1), P(4,1), A(5,1), M(6,1), D(7,1). Začnu tedy posouvat své objekty na tyto pozice. Nejdřív zkusím P.} \\
            }
        }
    }
    & & \\ \\

    \theutterance \stepcounter{utterance}  
    & & & \multicolumn{2}{p{0.3\linewidth}}{
        \cellcolor[rgb]{0.9,0.9,0.9}{
            \makecell[{{p{\linewidth}}}]{
                \texttt{\tiny{[GM$|$GM]}}
                \texttt{ŘEKNI: To je skvělý nápad! Jedna řada je velmi jednoduchá na koordinaci. Souhlasím s tvou konfigurací: S(1,1), O(2,1), T(3,1), P(4,1), A(5,1), M(6,1), D(7,1). Začnu tedy posouvat své objekty na tyto pozice. Nejdřív zkusím P.} \\
            }
        }
    }
    & & \\ \\

    \theutterance \stepcounter{utterance}  
    & & \multicolumn{4}{p{0.6\linewidth}}{
        \cellcolor[rgb]{0.9,0.9,0.9}{
            \makecell[{{p{\linewidth}}}]{
                \texttt{\tiny{[GM$\rangle$P2]}}
                \texttt{Vaše zpráva byla předána druhému hráči.} \\
\\ 
\\ 
\texttt{Právě hrajete kolo 2 z maximálně 28.} \\
\texttt{Dohromady jste nasbírali 0 z 16 pokut.} \\
\texttt{Druhý hráč poslal tuto zprávu:} \\
\\ 
\\ 
\texttt{"To je skvělý nápad! Jedna řada je velmi jednoduchá na koordinaci. Souhlasím s tvou konfigurací: S(1,1), O(2,1), T(3,1), P(4,1), A(5,1), M(6,1), D(7,1). Začnu tedy posouvat své objekty na tyto pozice. Nejdřív zkusím P."} \\
\\ 
\\ 
\texttt{Jaký je váš další příkaz?} \\
            }
        }
    }
    & \\ \\

    \theutterance \stepcounter{utterance}  
    & & & \multicolumn{4}{p{0.6\linewidth}}{
        \cellcolor[rgb]{0.9,0.9,0.9}{
            \makecell[{{p{\linewidth}}}]{
                \texttt{\tiny{[GM$\langle$P2]}}
                \texttt{POHYB: S, (1, 1)} \\
            }
        }
    }
     \\ \\

    \theutterance \stepcounter{utterance}  
    & & & \multicolumn{2}{p{0.3\linewidth}}{
        \cellcolor[rgb]{0.9,0.9,0.9}{
            \makecell[{{p{\linewidth}}}]{
                \texttt{\tiny{[GM$|$GM]}}
                \texttt{POHYB: S, (1, 1)} \\
            }
        }
    }
    & & \\ \\

    \theutterance \stepcounter{utterance}  
    & & & \multicolumn{2}{p{0.3\linewidth}}{
        \cellcolor[rgb]{0.9,0.9,0.9}{
            \makecell[{{p{\linewidth}}}]{
                \texttt{\tiny{[GM$|$GM]}}
                \texttt{S moved to (1, 1)} \\
            }
        }
    }
    & & \\ \\

    \theutterance \stepcounter{utterance}  
    & & \multicolumn{4}{p{0.6\linewidth}}{
        \cellcolor[rgb]{0.9,0.9,0.9}{
            \makecell[{{p{\linewidth}}}]{
                \texttt{\tiny{[P1$\langle$GM]}}
                \texttt{Vaše zpráva byla předána druhému hráči.} \\
\\ 
\\ 
\texttt{Právě hrajete kolo 2 z maximálně 28.} \\
\texttt{Dohromady jste nasbírali 0 z 16 pokut.} \\
\texttt{Druhý hráč přesunul objekt na své mřížce.} \\
\\ 
\\ 
\texttt{Jaký je váš další příkaz?} \\
            }
        }
    }
    & \\ \\

    \theutterance \stepcounter{utterance}  
    & \multicolumn{4}{p{0.6\linewidth}}{
        \cellcolor[rgb]{0.9,0.9,0.9}{
            \makecell[{{p{\linewidth}}}]{
                \texttt{\tiny{[P1$\rangle$GM]}}
                \texttt{POHYB: P, (4, 1)} \\
            }
        }
    }
    & & \\ \\

    \theutterance \stepcounter{utterance}  
    & & & \multicolumn{2}{p{0.3\linewidth}}{
        \cellcolor[rgb]{0.9,0.9,0.9}{
            \makecell[{{p{\linewidth}}}]{
                \texttt{\tiny{[GM$|$GM]}}
                \texttt{POHYB: P, (4, 1)} \\
            }
        }
    }
    & & \\ \\

    \theutterance \stepcounter{utterance}  
    & & & \multicolumn{2}{p{0.3\linewidth}}{
        \cellcolor[rgb]{0.9,0.9,0.9}{
            \makecell[{{p{\linewidth}}}]{
                \texttt{\tiny{[GM$|$GM]}}
                \texttt{P moved to (4, 1)} \\
            }
        }
    }
    & & \\ \\

    \theutterance \stepcounter{utterance}  
    & & \multicolumn{4}{p{0.6\linewidth}}{
        \cellcolor[rgb]{0.9,0.9,0.9}{
            \makecell[{{p{\linewidth}}}]{
                \texttt{\tiny{[GM$\rangle$P2]}}
                \texttt{Úspěšně přesunuto 'S' do (1,1). Vaše aktualizovaná mřížka vypadá takto:} \\
\texttt{```} \\
\texttt{ 1234567} \\
\texttt{╔═══════╗} \\
\texttt{║S◌◌◌◌◌◌║ 1} \\
\texttt{║◌◌◌◌O◌◌║ 2} \\
\texttt{║◌◌◌M◌┌─╢ 3} \\
\texttt{║◌◌A┌─┤◌║ 4} \\
\texttt{║◌◌◌│◌├─╢ 5} \\
\texttt{╟───┼─┘◌║ 6} \\
\texttt{║D◌T│P◌◌║ 7} \\
\texttt{╚═══╧═══╝} \\
\texttt{```} \\
\texttt{Právě hrajete kolo 3 z maximálně 28.} \\
\texttt{Dohromady jste nasbírali 0 z 16 pokut.} \\
\texttt{Druhý hráč přesunul objekt na své mřížce.} \\
\\ 
\\ 
\texttt{Jaký je váš další příkaz?} \\
            }
        }
    }
    & \\ \\

    \theutterance \stepcounter{utterance}  
    & & & \multicolumn{4}{p{0.6\linewidth}}{
        \cellcolor[rgb]{0.9,0.9,0.9}{
            \makecell[{{p{\linewidth}}}]{
                \texttt{\tiny{[GM$\langle$P2]}}
                \texttt{POHYB: O, (2, 1)} \\
            }
        }
    }
     \\ \\

    \theutterance \stepcounter{utterance}  
    & & & \multicolumn{2}{p{0.3\linewidth}}{
        \cellcolor[rgb]{0.9,0.9,0.9}{
            \makecell[{{p{\linewidth}}}]{
                \texttt{\tiny{[GM$|$GM]}}
                \texttt{POHYB: O, (2, 1)} \\
            }
        }
    }
    & & \\ \\

    \theutterance \stepcounter{utterance}  
    & & & \multicolumn{2}{p{0.3\linewidth}}{
        \cellcolor[rgb]{0.9,0.9,0.9}{
            \makecell[{{p{\linewidth}}}]{
                \texttt{\tiny{[GM$|$GM]}}
                \texttt{O moved to (2, 1)} \\
            }
        }
    }
    & & \\ \\

    \theutterance \stepcounter{utterance}  
    & & \multicolumn{4}{p{0.6\linewidth}}{
        \cellcolor[rgb]{0.9,0.9,0.9}{
            \makecell[{{p{\linewidth}}}]{
                \texttt{\tiny{[P1$\langle$GM]}}
                \texttt{Úspěšně přesunuto 'P' do (4,1). Vaše aktualizovaná mřížka vypadá takto:} \\
\texttt{```} \\
\texttt{ 1234567} \\
\texttt{╔═══════╗} \\
\texttt{║◌◌◌P◌T◌║ 1} \\
\texttt{║◌◌◌◌◌◌◌║ 2} \\
\texttt{║◌◌AD◌┌─╢ 3} \\
\texttt{║◌◌◌┌─┤S║ 4} \\
\texttt{║O◌◌│◌├─╢ 5} \\
\texttt{╟───┼─┘◌║ 6} \\
\texttt{║◌M◌│◌◌◌║ 7} \\
\texttt{╚═══╧═══╝} \\
\texttt{```} \\
\texttt{Právě hrajete kolo 3 z maximálně 28.} \\
\texttt{Dohromady jste nasbírali 0 z 16 pokut.} \\
\texttt{Druhý hráč přesunul objekt na své mřížce.} \\
\\ 
\\ 
\texttt{Jaký je váš další příkaz?} \\
            }
        }
    }
    & \\ \\

    \theutterance \stepcounter{utterance}  
    & \multicolumn{4}{p{0.6\linewidth}}{
        \cellcolor[rgb]{0.9,0.9,0.9}{
            \makecell[{{p{\linewidth}}}]{
                \texttt{\tiny{[P1$\rangle$GM]}}
                \texttt{POHYB: T, (3, 1)} \\
            }
        }
    }
    & & \\ \\

    \theutterance \stepcounter{utterance}  
    & & & \multicolumn{2}{p{0.3\linewidth}}{
        \cellcolor[rgb]{0.9,0.9,0.9}{
            \makecell[{{p{\linewidth}}}]{
                \texttt{\tiny{[GM$|$GM]}}
                \texttt{POHYB: T, (3, 1)} \\
            }
        }
    }
    & & \\ \\

    \theutterance \stepcounter{utterance}  
    & & & \multicolumn{2}{p{0.3\linewidth}}{
        \cellcolor[rgb]{0.9,0.9,0.9}{
            \makecell[{{p{\linewidth}}}]{
                \texttt{\tiny{[GM$|$GM]}}
                \texttt{T moved to (3, 1)} \\
            }
        }
    }
    & & \\ \\

    \theutterance \stepcounter{utterance}  
    & & \multicolumn{4}{p{0.6\linewidth}}{
        \cellcolor[rgb]{0.9,0.9,0.9}{
            \makecell[{{p{\linewidth}}}]{
                \texttt{\tiny{[GM$\rangle$P2]}}
                \texttt{Úspěšně přesunuto 'O' do (2,1). Vaše aktualizovaná mřížka vypadá takto:} \\
\texttt{```} \\
\texttt{ 1234567} \\
\texttt{╔═══════╗} \\
\texttt{║SO◌◌◌◌◌║ 1} \\
\texttt{║◌◌◌◌◌◌◌║ 2} \\
\texttt{║◌◌◌M◌┌─╢ 3} \\
\texttt{║◌◌A┌─┤◌║ 4} \\
\texttt{║◌◌◌│◌├─╢ 5} \\
\texttt{╟───┼─┘◌║ 6} \\
\texttt{║D◌T│P◌◌║ 7} \\
\texttt{╚═══╧═══╝} \\
\texttt{```} \\
\texttt{Právě hrajete kolo 4 z maximálně 28.} \\
\texttt{Dohromady jste nasbírali 0 z 16 pokut.} \\
\texttt{Druhý hráč přesunul objekt na své mřížce.} \\
\\ 
\\ 
\texttt{Jaký je váš další příkaz?} \\
            }
        }
    }
    & \\ \\

    \theutterance \stepcounter{utterance}  
    & & & \multicolumn{4}{p{0.6\linewidth}}{
        \cellcolor[rgb]{0.9,0.9,0.9}{
            \makecell[{{p{\linewidth}}}]{
                \texttt{\tiny{[GM$\langle$P2]}}
                \texttt{POHYB: T, (3, 1)} \\
            }
        }
    }
     \\ \\

    \theutterance \stepcounter{utterance}  
    & & & \multicolumn{2}{p{0.3\linewidth}}{
        \cellcolor[rgb]{0.9,0.9,0.9}{
            \makecell[{{p{\linewidth}}}]{
                \texttt{\tiny{[GM$|$GM]}}
                \texttt{POHYB: T, (3, 1)} \\
            }
        }
    }
    & & \\ \\

    \theutterance \stepcounter{utterance}  
    & & & \multicolumn{2}{p{0.3\linewidth}}{
        \cellcolor[rgb]{0.9,0.9,0.9}{
            \makecell[{{p{\linewidth}}}]{
                \texttt{\tiny{[GM$|$GM]}}
                \texttt{T moved to (3, 1)} \\
            }
        }
    }
    & & \\ \\

    \theutterance \stepcounter{utterance}  
    & & \multicolumn{4}{p{0.6\linewidth}}{
        \cellcolor[rgb]{0.9,0.9,0.9}{
            \makecell[{{p{\linewidth}}}]{
                \texttt{\tiny{[P1$\langle$GM]}}
                \texttt{Úspěšně přesunuto 'T' do (3,1). Vaše aktualizovaná mřížka vypadá takto:} \\
\texttt{```} \\
\texttt{ 1234567} \\
\texttt{╔═══════╗} \\
\texttt{║◌◌TP◌◌◌║ 1} \\
\texttt{║◌◌◌◌◌◌◌║ 2} \\
\texttt{║◌◌AD◌┌─╢ 3} \\
\texttt{║◌◌◌┌─┤S║ 4} \\
\texttt{║O◌◌│◌├─╢ 5} \\
\texttt{╟───┼─┘◌║ 6} \\
\texttt{║◌M◌│◌◌◌║ 7} \\
\texttt{╚═══╧═══╝} \\
\texttt{```} \\
\texttt{Právě hrajete kolo 4 z maximálně 28.} \\
\texttt{Dohromady jste nasbírali 0 z 16 pokut.} \\
\texttt{Druhý hráč přesunul objekt na své mřížce.} \\
\\ 
\\ 
\texttt{Jaký je váš další příkaz?} \\
            }
        }
    }
    & \\ \\

    \theutterance \stepcounter{utterance}  
    & \multicolumn{4}{p{0.6\linewidth}}{
        \cellcolor[rgb]{0.9,0.9,0.9}{
            \makecell[{{p{\linewidth}}}]{
                \texttt{\tiny{[P1$\rangle$GM]}}
                \texttt{POHYB: S, (1, 1)} \\
            }
        }
    }
    & & \\ \\

    \theutterance \stepcounter{utterance}  
    & & & \multicolumn{2}{p{0.3\linewidth}}{
        \cellcolor[rgb]{0.9,0.9,0.9}{
            \makecell[{{p{\linewidth}}}]{
                \texttt{\tiny{[GM$|$GM]}}
                \texttt{POHYB: S, (1, 1)} \\
            }
        }
    }
    & & \\ \\

    \theutterance \stepcounter{utterance}  
    & & & \multicolumn{2}{p{0.3\linewidth}}{
        \cellcolor[rgb]{0.9,0.9,0.9}{
            \makecell[{{p{\linewidth}}}]{
                \texttt{\tiny{[GM$|$GM]}}
                \texttt{S moved to (1, 1)} \\
            }
        }
    }
    & & \\ \\

    \theutterance \stepcounter{utterance}  
    & & \multicolumn{4}{p{0.6\linewidth}}{
        \cellcolor[rgb]{0.9,0.9,0.9}{
            \makecell[{{p{\linewidth}}}]{
                \texttt{\tiny{[GM$\rangle$P2]}}
                \texttt{Úspěšně přesunuto 'T' do (3,1). Vaše aktualizovaná mřížka vypadá takto:} \\
\texttt{```} \\
\texttt{ 1234567} \\
\texttt{╔═══════╗} \\
\texttt{║SOT◌◌◌◌║ 1} \\
\texttt{║◌◌◌◌◌◌◌║ 2} \\
\texttt{║◌◌◌M◌┌─╢ 3} \\
\texttt{║◌◌A┌─┤◌║ 4} \\
\texttt{║◌◌◌│◌├─╢ 5} \\
\texttt{╟───┼─┘◌║ 6} \\
\texttt{║D◌◌│P◌◌║ 7} \\
\texttt{╚═══╧═══╝} \\
\texttt{```} \\
\texttt{Právě hrajete kolo 5 z maximálně 28.} \\
\texttt{Dohromady jste nasbírali 0 z 16 pokut.} \\
\texttt{Druhý hráč přesunul objekt na své mřížce.} \\
\\ 
\\ 
\texttt{Jaký je váš další příkaz?} \\
            }
        }
    }
    & \\ \\

    \theutterance \stepcounter{utterance}  
    & & & \multicolumn{4}{p{0.6\linewidth}}{
        \cellcolor[rgb]{0.9,0.9,0.9}{
            \makecell[{{p{\linewidth}}}]{
                \texttt{\tiny{[GM$\langle$P2]}}
                \texttt{POHYB: P, (4, 1)} \\
            }
        }
    }
     \\ \\

    \theutterance \stepcounter{utterance}  
    & & & \multicolumn{2}{p{0.3\linewidth}}{
        \cellcolor[rgb]{0.9,0.9,0.9}{
            \makecell[{{p{\linewidth}}}]{
                \texttt{\tiny{[GM$|$GM]}}
                \texttt{POHYB: P, (4, 1)} \\
            }
        }
    }
    & & \\ \\

    \theutterance \stepcounter{utterance}  
    & & & \multicolumn{2}{p{0.3\linewidth}}{
        \cellcolor[rgb]{0.9,0.9,0.9}{
            \makecell[{{p{\linewidth}}}]{
                \texttt{\tiny{[GM$|$GM]}}
                \texttt{P moved to (4, 1)} \\
            }
        }
    }
    & & \\ \\

    \theutterance \stepcounter{utterance}  
    & & \multicolumn{4}{p{0.6\linewidth}}{
        \cellcolor[rgb]{0.9,0.9,0.9}{
            \makecell[{{p{\linewidth}}}]{
                \texttt{\tiny{[P1$\langle$GM]}}
                \texttt{Úspěšně přesunuto 'S' do (1,1). Vaše aktualizovaná mřížka vypadá takto:} \\
\texttt{```} \\
\texttt{ 1234567} \\
\texttt{╔═══════╗} \\
\texttt{║S◌TP◌◌◌║ 1} \\
\texttt{║◌◌◌◌◌◌◌║ 2} \\
\texttt{║◌◌AD◌┌─╢ 3} \\
\texttt{║◌◌◌┌─┤◌║ 4} \\
\texttt{║O◌◌│◌├─╢ 5} \\
\texttt{╟───┼─┘◌║ 6} \\
\texttt{║◌M◌│◌◌◌║ 7} \\
\texttt{╚═══╧═══╝} \\
\texttt{```} \\
\texttt{Právě hrajete kolo 5 z maximálně 28.} \\
\texttt{Dohromady jste nasbírali 0 z 16 pokut.} \\
\texttt{Druhý hráč přesunul objekt na své mřížce.} \\
\\ 
\\ 
\texttt{Jaký je váš další příkaz?} \\
            }
        }
    }
    & \\ \\

    \theutterance \stepcounter{utterance}  
    & \multicolumn{4}{p{0.6\linewidth}}{
        \cellcolor[rgb]{0.9,0.9,0.9}{
            \makecell[{{p{\linewidth}}}]{
                \texttt{\tiny{[P1$\rangle$GM]}}
                \texttt{POHYB: O, (2, 1)} \\
            }
        }
    }
    & & \\ \\

    \theutterance \stepcounter{utterance}  
    & & & \multicolumn{2}{p{0.3\linewidth}}{
        \cellcolor[rgb]{0.9,0.9,0.9}{
            \makecell[{{p{\linewidth}}}]{
                \texttt{\tiny{[GM$|$GM]}}
                \texttt{POHYB: O, (2, 1)} \\
            }
        }
    }
    & & \\ \\

    \theutterance \stepcounter{utterance}  
    & & & \multicolumn{2}{p{0.3\linewidth}}{
        \cellcolor[rgb]{0.9,0.9,0.9}{
            \makecell[{{p{\linewidth}}}]{
                \texttt{\tiny{[GM$|$GM]}}
                \texttt{O moved to (2, 1)} \\
            }
        }
    }
    & & \\ \\

    \theutterance \stepcounter{utterance}  
    & & \multicolumn{4}{p{0.6\linewidth}}{
        \cellcolor[rgb]{0.9,0.9,0.9}{
            \makecell[{{p{\linewidth}}}]{
                \texttt{\tiny{[GM$\rangle$P2]}}
                \texttt{Úspěšně přesunuto 'P' do (4,1). Vaše aktualizovaná mřížka vypadá takto:} \\
\texttt{```} \\
\texttt{ 1234567} \\
\texttt{╔═══════╗} \\
\texttt{║SOTP◌◌◌║ 1} \\
\texttt{║◌◌◌◌◌◌◌║ 2} \\
\texttt{║◌◌◌M◌┌─╢ 3} \\
\texttt{║◌◌A┌─┤◌║ 4} \\
\texttt{║◌◌◌│◌├─╢ 5} \\
\texttt{╟───┼─┘◌║ 6} \\
\texttt{║D◌◌│◌◌◌║ 7} \\
\texttt{╚═══╧═══╝} \\
\texttt{```} \\
\texttt{Právě hrajete kolo 6 z maximálně 28.} \\
\texttt{Dohromady jste nasbírali 0 z 16 pokut.} \\
\texttt{Druhý hráč přesunul objekt na své mřížce.} \\
\\ 
\\ 
\texttt{Jaký je váš další příkaz?} \\
            }
        }
    }
    & \\ \\

    \theutterance \stepcounter{utterance}  
    & & & \multicolumn{4}{p{0.6\linewidth}}{
        \cellcolor[rgb]{0.9,0.9,0.9}{
            \makecell[{{p{\linewidth}}}]{
                \texttt{\tiny{[GM$\langle$P2]}}
                \texttt{POHYB: A, (5, 1)} \\
            }
        }
    }
     \\ \\

    \theutterance \stepcounter{utterance}  
    & & & \multicolumn{2}{p{0.3\linewidth}}{
        \cellcolor[rgb]{0.9,0.9,0.9}{
            \makecell[{{p{\linewidth}}}]{
                \texttt{\tiny{[GM$|$GM]}}
                \texttt{POHYB: A, (5, 1)} \\
            }
        }
    }
    & & \\ \\

    \theutterance \stepcounter{utterance}  
    & & & \multicolumn{2}{p{0.3\linewidth}}{
        \cellcolor[rgb]{0.9,0.9,0.9}{
            \makecell[{{p{\linewidth}}}]{
                \texttt{\tiny{[GM$|$GM]}}
                \texttt{A moved to (5, 1)} \\
            }
        }
    }
    & & \\ \\

    \theutterance \stepcounter{utterance}  
    & & \multicolumn{4}{p{0.6\linewidth}}{
        \cellcolor[rgb]{0.9,0.9,0.9}{
            \makecell[{{p{\linewidth}}}]{
                \texttt{\tiny{[P1$\langle$GM]}}
                \texttt{Úspěšně přesunuto 'O' do (2,1). Vaše aktualizovaná mřížka vypadá takto:} \\
\texttt{```} \\
\texttt{ 1234567} \\
\texttt{╔═══════╗} \\
\texttt{║SOTP◌◌◌║ 1} \\
\texttt{║◌◌◌◌◌◌◌║ 2} \\
\texttt{║◌◌AD◌┌─╢ 3} \\
\texttt{║◌◌◌┌─┤◌║ 4} \\
\texttt{║◌◌◌│◌├─╢ 5} \\
\texttt{╟───┼─┘◌║ 6} \\
\texttt{║◌M◌│◌◌◌║ 7} \\
\texttt{╚═══╧═══╝} \\
\texttt{```} \\
\texttt{Právě hrajete kolo 6 z maximálně 28.} \\
\texttt{Dohromady jste nasbírali 0 z 16 pokut.} \\
\texttt{Druhý hráč přesunul objekt na své mřížce.} \\
\\ 
\\ 
\texttt{Jaký je váš další příkaz?} \\
            }
        }
    }
    & \\ \\

    \theutterance \stepcounter{utterance}  
    & \multicolumn{4}{p{0.6\linewidth}}{
        \cellcolor[rgb]{0.9,0.9,0.9}{
            \makecell[{{p{\linewidth}}}]{
                \texttt{\tiny{[P1$\rangle$GM]}}
                \texttt{POHYB: A, (5, 1)} \\
            }
        }
    }
    & & \\ \\

    \theutterance \stepcounter{utterance}  
    & & & \multicolumn{2}{p{0.3\linewidth}}{
        \cellcolor[rgb]{0.9,0.9,0.9}{
            \makecell[{{p{\linewidth}}}]{
                \texttt{\tiny{[GM$|$GM]}}
                \texttt{POHYB: A, (5, 1)} \\
            }
        }
    }
    & & \\ \\

    \theutterance \stepcounter{utterance}  
    & & & \multicolumn{2}{p{0.3\linewidth}}{
        \cellcolor[rgb]{0.9,0.9,0.9}{
            \makecell[{{p{\linewidth}}}]{
                \texttt{\tiny{[GM$|$GM]}}
                \texttt{A moved to (5, 1)} \\
            }
        }
    }
    & & \\ \\

    \theutterance \stepcounter{utterance}  
    & & \multicolumn{4}{p{0.6\linewidth}}{
        \cellcolor[rgb]{0.9,0.9,0.9}{
            \makecell[{{p{\linewidth}}}]{
                \texttt{\tiny{[GM$\rangle$P2]}}
                \texttt{Úspěšně přesunuto 'A' do (5,1). Vaše aktualizovaná mřížka vypadá takto:} \\
\texttt{```} \\
\texttt{ 1234567} \\
\texttt{╔═══════╗} \\
\texttt{║SOTPA◌◌║ 1} \\
\texttt{║◌◌◌◌◌◌◌║ 2} \\
\texttt{║◌◌◌M◌┌─╢ 3} \\
\texttt{║◌◌◌┌─┤◌║ 4} \\
\texttt{║◌◌◌│◌├─╢ 5} \\
\texttt{╟───┼─┘◌║ 6} \\
\texttt{║D◌◌│◌◌◌║ 7} \\
\texttt{╚═══╧═══╝} \\
\texttt{```} \\
\texttt{Právě hrajete kolo 7 z maximálně 28.} \\
\texttt{Dohromady jste nasbírali 0 z 16 pokut.} \\
\texttt{Druhý hráč přesunul objekt na své mřížce.} \\
\\ 
\\ 
\texttt{Jaký je váš další příkaz?} \\
            }
        }
    }
    & \\ \\

    \theutterance \stepcounter{utterance}  
    & & & \multicolumn{4}{p{0.6\linewidth}}{
        \cellcolor[rgb]{0.9,0.9,0.9}{
            \makecell[{{p{\linewidth}}}]{
                \texttt{\tiny{[GM$\langle$P2]}}
                \texttt{POHYB: M, (6, 1)} \\
            }
        }
    }
     \\ \\

    \theutterance \stepcounter{utterance}  
    & & & \multicolumn{2}{p{0.3\linewidth}}{
        \cellcolor[rgb]{0.9,0.9,0.9}{
            \makecell[{{p{\linewidth}}}]{
                \texttt{\tiny{[GM$|$GM]}}
                \texttt{POHYB: M, (6, 1)} \\
            }
        }
    }
    & & \\ \\

    \theutterance \stepcounter{utterance}  
    & & & \multicolumn{2}{p{0.3\linewidth}}{
        \cellcolor[rgb]{0.9,0.9,0.9}{
            \makecell[{{p{\linewidth}}}]{
                \texttt{\tiny{[GM$|$GM]}}
                \texttt{M moved to (6, 1)} \\
            }
        }
    }
    & & \\ \\

    \theutterance \stepcounter{utterance}  
    & & \multicolumn{4}{p{0.6\linewidth}}{
        \cellcolor[rgb]{0.9,0.9,0.9}{
            \makecell[{{p{\linewidth}}}]{
                \texttt{\tiny{[P1$\langle$GM]}}
                \texttt{Úspěšně přesunuto 'A' do (5,1). Vaše aktualizovaná mřížka vypadá takto:} \\
\texttt{```} \\
\texttt{ 1234567} \\
\texttt{╔═══════╗} \\
\texttt{║SOTPA◌◌║ 1} \\
\texttt{║◌◌◌◌◌◌◌║ 2} \\
\texttt{║◌◌◌D◌┌─╢ 3} \\
\texttt{║◌◌◌┌─┤◌║ 4} \\
\texttt{║◌◌◌│◌├─╢ 5} \\
\texttt{╟───┼─┘◌║ 6} \\
\texttt{║◌M◌│◌◌◌║ 7} \\
\texttt{╚═══╧═══╝} \\
\texttt{```} \\
\texttt{Právě hrajete kolo 7 z maximálně 28.} \\
\texttt{Dohromady jste nasbírali 0 z 16 pokut.} \\
\texttt{Druhý hráč přesunul objekt na své mřížce.} \\
\\ 
\\ 
\texttt{Jaký je váš další příkaz?} \\
            }
        }
    }
    & \\ \\

    \theutterance \stepcounter{utterance}  
    & \multicolumn{4}{p{0.6\linewidth}}{
        \cellcolor[rgb]{0.9,0.9,0.9}{
            \makecell[{{p{\linewidth}}}]{
                \texttt{\tiny{[P1$\rangle$GM]}}
                \texttt{POHYB: M, (6, 1)} \\
            }
        }
    }
    & & \\ \\

    \theutterance \stepcounter{utterance}  
    & & & \multicolumn{2}{p{0.3\linewidth}}{
        \cellcolor[rgb]{0.9,0.9,0.9}{
            \makecell[{{p{\linewidth}}}]{
                \texttt{\tiny{[GM$|$GM]}}
                \texttt{POHYB: M, (6, 1)} \\
            }
        }
    }
    & & \\ \\

    \theutterance \stepcounter{utterance}  
    & & & \multicolumn{2}{p{0.3\linewidth}}{
        \cellcolor[rgb]{0.9,0.9,0.9}{
            \makecell[{{p{\linewidth}}}]{
                \texttt{\tiny{[GM$|$GM]}}
                \texttt{M moved to (6, 1)} \\
            }
        }
    }
    & & \\ \\

    \theutterance \stepcounter{utterance}  
    & & \multicolumn{4}{p{0.6\linewidth}}{
        \cellcolor[rgb]{0.9,0.9,0.9}{
            \makecell[{{p{\linewidth}}}]{
                \texttt{\tiny{[GM$\rangle$P2]}}
                \texttt{Úspěšně přesunuto 'M' do (6,1). Vaše aktualizovaná mřížka vypadá takto:} \\
\texttt{```} \\
\texttt{ 1234567} \\
\texttt{╔═══════╗} \\
\texttt{║SOTPAM◌║ 1} \\
\texttt{║◌◌◌◌◌◌◌║ 2} \\
\texttt{║◌◌◌◌◌┌─╢ 3} \\
\texttt{║◌◌◌┌─┤◌║ 4} \\
\texttt{║◌◌◌│◌├─╢ 5} \\
\texttt{╟───┼─┘◌║ 6} \\
\texttt{║D◌◌│◌◌◌║ 7} \\
\texttt{╚═══╧═══╝} \\
\texttt{```} \\
\texttt{Právě hrajete kolo 8 z maximálně 28.} \\
\texttt{Dohromady jste nasbírali 0 z 16 pokut.} \\
\texttt{Druhý hráč přesunul objekt na své mřížce.} \\
\\ 
\\ 
\texttt{Jaký je váš další příkaz?} \\
            }
        }
    }
    & \\ \\

    \theutterance \stepcounter{utterance}  
    & & & \multicolumn{4}{p{0.6\linewidth}}{
        \cellcolor[rgb]{0.9,0.9,0.9}{
            \makecell[{{p{\linewidth}}}]{
                \texttt{\tiny{[GM$\langle$P2]}}
                \texttt{POHYB: D, (7, 1)} \\
            }
        }
    }
     \\ \\

    \theutterance \stepcounter{utterance}  
    & & & \multicolumn{2}{p{0.3\linewidth}}{
        \cellcolor[rgb]{0.9,0.9,0.9}{
            \makecell[{{p{\linewidth}}}]{
                \texttt{\tiny{[GM$|$GM]}}
                \texttt{POHYB: D, (7, 1)} \\
            }
        }
    }
    & & \\ \\

    \theutterance \stepcounter{utterance}  
    & & & \multicolumn{2}{p{0.3\linewidth}}{
        \cellcolor[rgb]{0.9,0.9,0.9}{
            \makecell[{{p{\linewidth}}}]{
                \texttt{\tiny{[GM$|$GM]}}
                \texttt{D moved to (7, 1)} \\
            }
        }
    }
    & & \\ \\

    \theutterance \stepcounter{utterance}  
    & & \multicolumn{4}{p{0.6\linewidth}}{
        \cellcolor[rgb]{0.9,0.9,0.9}{
            \makecell[{{p{\linewidth}}}]{
                \texttt{\tiny{[P1$\langle$GM]}}
                \texttt{Úspěšně přesunuto 'M' do (6,1). Vaše aktualizovaná mřížka vypadá takto:} \\
\texttt{```} \\
\texttt{ 1234567} \\
\texttt{╔═══════╗} \\
\texttt{║SOTPAM◌║ 1} \\
\texttt{║◌◌◌◌◌◌◌║ 2} \\
\texttt{║◌◌◌D◌┌─╢ 3} \\
\texttt{║◌◌◌┌─┤◌║ 4} \\
\texttt{║◌◌◌│◌├─╢ 5} \\
\texttt{╟───┼─┘◌║ 6} \\
\texttt{║◌◌◌│◌◌◌║ 7} \\
\texttt{╚═══╧═══╝} \\
\texttt{```} \\
\texttt{Právě hrajete kolo 8 z maximálně 28.} \\
\texttt{Dohromady jste nasbírali 0 z 16 pokut.} \\
\texttt{Druhý hráč přesunul objekt na své mřížce.} \\
\\ 
\\ 
\texttt{Jaký je váš další příkaz?} \\
            }
        }
    }
    & \\ \\

    \theutterance \stepcounter{utterance}  
    & \multicolumn{4}{p{0.6\linewidth}}{
        \cellcolor[rgb]{0.9,0.9,0.9}{
            \makecell[{{p{\linewidth}}}]{
                \texttt{\tiny{[P1$\rangle$GM]}}
                \texttt{POHYB: D, (7, 1)} \\
            }
        }
    }
    & & \\ \\

    \theutterance \stepcounter{utterance}  
    & & & \multicolumn{2}{p{0.3\linewidth}}{
        \cellcolor[rgb]{0.9,0.9,0.9}{
            \makecell[{{p{\linewidth}}}]{
                \texttt{\tiny{[GM$|$GM]}}
                \texttt{POHYB: D, (7, 1)} \\
            }
        }
    }
    & & \\ \\

    \theutterance \stepcounter{utterance}  
    & & & \multicolumn{2}{p{0.3\linewidth}}{
        \cellcolor[rgb]{0.9,0.9,0.9}{
            \makecell[{{p{\linewidth}}}]{
                \texttt{\tiny{[GM$|$GM]}}
                \texttt{D moved to (7, 1)} \\
            }
        }
    }
    & & \\ \\

    \theutterance \stepcounter{utterance}  
    & & \multicolumn{4}{p{0.6\linewidth}}{
        \cellcolor[rgb]{0.9,0.9,0.9}{
            \makecell[{{p{\linewidth}}}]{
                \texttt{\tiny{[GM$\rangle$P2]}}
                \texttt{Úspěšně přesunuto 'D' do (7,1). Vaše aktualizovaná mřížka vypadá takto:} \\
\texttt{```} \\
\texttt{ 1234567} \\
\texttt{╔═══════╗} \\
\texttt{║SOTPAMD║ 1} \\
\texttt{║◌◌◌◌◌◌◌║ 2} \\
\texttt{║◌◌◌◌◌┌─╢ 3} \\
\texttt{║◌◌◌┌─┤◌║ 4} \\
\texttt{║◌◌◌│◌├─╢ 5} \\
\texttt{╟───┼─┘◌║ 6} \\
\texttt{║◌◌◌│◌◌◌║ 7} \\
\texttt{╚═══╧═══╝} \\
\texttt{```} \\
\texttt{Právě hrajete kolo 9 z maximálně 28.} \\
\texttt{Dohromady jste nasbírali 0 z 16 pokut.} \\
\texttt{Druhý hráč přesunul objekt na své mřížce.} \\
\\ 
\\ 
\texttt{Jaký je váš další příkaz?} \\
            }
        }
    }
    & \\ \\

    \theutterance \stepcounter{utterance}  
    & & & \multicolumn{4}{p{0.6\linewidth}}{
        \cellcolor[rgb]{0.9,0.9,0.9}{
            \makecell[{{p{\linewidth}}}]{
                \texttt{\tiny{[GM$\langle$P2]}}
                \texttt{ŘEKNI: hotovo?} \\
            }
        }
    }
     \\ \\

    \theutterance \stepcounter{utterance}  
    & & & \multicolumn{2}{p{0.3\linewidth}}{
        \cellcolor[rgb]{0.9,0.9,0.9}{
            \makecell[{{p{\linewidth}}}]{
                \texttt{\tiny{[GM$|$GM]}}
                \texttt{ŘEKNI: hotovo?} \\
            }
        }
    }
    & & \\ \\

    \theutterance \stepcounter{utterance}  
    & & \multicolumn{4}{p{0.6\linewidth}}{
        \cellcolor[rgb]{0.9,0.9,0.9}{
            \makecell[{{p{\linewidth}}}]{
                \texttt{\tiny{[P1$\langle$GM]}}
                \texttt{Úspěšně přesunuto 'D' do (7,1). Vaše aktualizovaná mřížka vypadá takto:} \\
\texttt{```} \\
\texttt{ 1234567} \\
\texttt{╔═══════╗} \\
\texttt{║SOTPAMD║ 1} \\
\texttt{║◌◌◌◌◌◌◌║ 2} \\
\texttt{║◌◌◌◌◌┌─╢ 3} \\
\texttt{║◌◌◌┌─┤◌║ 4} \\
\texttt{║◌◌◌│◌├─╢ 5} \\
\texttt{╟───┼─┘◌║ 6} \\
\texttt{║◌◌◌│◌◌◌║ 7} \\
\texttt{╚═══╧═══╝} \\
\texttt{```} \\
\texttt{Právě hrajete kolo 9 z maximálně 28.} \\
\texttt{Dohromady jste nasbírali 0 z 16 pokut.} \\
\texttt{Druhý hráč poslal tuto zprávu:} \\
\\ 
\\ 
\texttt{"hotovo?"} \\
\\ 
\\ 
\texttt{Jaký je váš další příkaz?} \\
            }
        }
    }
    & \\ \\

    \theutterance \stepcounter{utterance}  
    & \multicolumn{4}{p{0.6\linewidth}}{
        \cellcolor[rgb]{0.9,0.9,0.9}{
            \makecell[{{p{\linewidth}}}]{
                \texttt{\tiny{[P1$\rangle$GM]}}
                \texttt{ŘEKNI: hotovo!} \\
            }
        }
    }
    & & \\ \\

    \theutterance \stepcounter{utterance}  
    & & & \multicolumn{2}{p{0.3\linewidth}}{
        \cellcolor[rgb]{0.9,0.9,0.9}{
            \makecell[{{p{\linewidth}}}]{
                \texttt{\tiny{[GM$|$GM]}}
                \texttt{ŘEKNI: hotovo!} \\
            }
        }
    }
    & & \\ \\

    \theutterance \stepcounter{utterance}  
    & & & \multicolumn{2}{p{0.3\linewidth}}{
        \cellcolor[rgb]{0.9,0.9,0.9}{
            \makecell[{{p{\linewidth}}}]{
                \texttt{\tiny{[GM$|$GM]}}
                \texttt{true} \\
            }
        }
    }
    & & \\ \\

    \theutterance \stepcounter{utterance}  
    & & & \multicolumn{2}{p{0.3\linewidth}}{
        \cellcolor[rgb]{0.9,0.9,0.9}{
            \makecell[{{p{\linewidth}}}]{
                \texttt{\tiny{[GM$|$GM]}}
                \texttt{Initial states:} \\
\texttt{```} \\
\texttt{Player 1:} \\
\texttt{ 1234567} \\
\texttt{╔═══════╗} \\
\texttt{║◌◌◌◌◌T◌║ 1} \\
\texttt{║P◌◌◌◌◌◌║ 2} \\
\texttt{║◌◌AD◌┌─╢ 3} \\
\texttt{║◌◌◌┌─┤S║ 4} \\
\texttt{║O◌◌│◌├─╢ 5} \\
\texttt{╟───┼─┘◌║ 6} \\
\texttt{║◌M◌│◌◌◌║ 7} \\
\texttt{╚═══╧═══╝} \\
\\ 
\\ 
\texttt{Player 2:} \\
\texttt{ 1234567} \\
\texttt{╔═══════╗} \\
\texttt{║◌◌◌◌◌S◌║ 1} \\
\texttt{║◌◌◌◌O◌◌║ 2} \\
\texttt{║◌◌◌M◌┌─╢ 3} \\
\texttt{║◌◌A┌─┤◌║ 4} \\
\texttt{║◌◌◌│◌├─╢ 5} \\
\texttt{╟───┼─┘◌║ 6} \\
\texttt{║D◌T│P◌◌║ 7} \\
\texttt{╚═══╧═══╝} \\
\\ 
\texttt{```} \\
            }
        }
    }
    & & \\ \\

    \theutterance \stepcounter{utterance}  
    & & & \multicolumn{2}{p{0.3\linewidth}}{
        \cellcolor[rgb]{0.9,0.9,0.9}{
            \makecell[{{p{\linewidth}}}]{
                \texttt{\tiny{[GM$|$GM]}}
                \texttt{End states:} \\
\texttt{```} \\
\texttt{Player 1:} \\
\texttt{ 1234567} \\
\texttt{╔═══════╗} \\
\texttt{║SOTPAMD║ 1} \\
\texttt{║◌◌◌◌◌◌◌║ 2} \\
\texttt{║◌◌◌◌◌┌─╢ 3} \\
\texttt{║◌◌◌┌─┤◌║ 4} \\
\texttt{║◌◌◌│◌├─╢ 5} \\
\texttt{╟───┼─┘◌║ 6} \\
\texttt{║◌◌◌│◌◌◌║ 7} \\
\texttt{╚═══╧═══╝} \\
\\ 
\\ 
\texttt{Player 2:} \\
\texttt{ 1234567} \\
\texttt{╔═══════╗} \\
\texttt{║SOTPAMD║ 1} \\
\texttt{║◌◌◌◌◌◌◌║ 2} \\
\texttt{║◌◌◌◌◌┌─╢ 3} \\
\texttt{║◌◌◌┌─┤◌║ 4} \\
\texttt{║◌◌◌│◌├─╢ 5} \\
\texttt{╟───┼─┘◌║ 6} \\
\texttt{║◌◌◌│◌◌◌║ 7} \\
\texttt{╚═══╧═══╝} \\
\\ 
\texttt{```} \\
            }
        }
    }
    & & \\ \\

    \theutterance \stepcounter{utterance}  
    & & & \multicolumn{2}{p{0.3\linewidth}}{
        \cellcolor[rgb]{0.9,0.9,0.9}{
            \makecell[{{p{\linewidth}}}]{
                \texttt{\tiny{[GM$|$GM]}}
                \texttt{* success: True} \\
\texttt{* lose: False} \\
\texttt{* aborted: False} \\
\texttt{{-}{-}{-}{-}{-}{-}{-}} \\
\texttt{* Shifts: 9.00} \\
\texttt{* Max Shifts: 12.00} \\
\texttt{* Min Shifts: 6.00} \\
\texttt{* End Distance Sum: 0.00} \\
\texttt{* Init Distance Sum: 31.75} \\
\texttt{* Expected Distance Sum: 29.33} \\
\texttt{* Penalties: 0.00} \\
\texttt{* Max Penalties: 16.00} \\
\texttt{* Rounds: 9.00} \\
\texttt{* Max Rounds: 28.00} \\
\texttt{* Object Count: 7.00} \\
            }
        }
    }
    & & \\ \\

\end{supertabular}
}

\end{document}
